\documentclass[12pt,a4paper]{article}

\usepackage[utf8]{inputenc}
\usepackage[T1]{fontenc}
\usepackage{lmodern}
\usepackage{charter}
\usepackage[ngerman]{babel}
\usepackage[margin=1.6cm]{geometry}

\usepackage{enumitem}
\usepackage{titlesec}
\usepackage[document]{ragged2e}
\usepackage[shortcuts]{extdash}
\usepackage{booktabs}

\hyphenpenalty=10000
\exhyphenpenalty=0
\emergencystretch=2em

\pagestyle{empty}

\setlist[itemize]{
  noitemsep,
  topsep=5pt,
  leftmargin=2em,
}

\setlist[description]{
  leftmargin=!,
  labelwidth=1.5cm,
  style=nextline
}

\titleformat{\section}{\fontsize{18}{22}\bfseries}{}{0em}{}[{\titlerule[1.5pt]}]
\titlespacing*{\section}{0pt}{18pt}{10pt}

\begin{document}

\noindent
{\fontsize{24}{28}\selectfont\bfseries Studienarbeit} \\
\begin{minipage}[t]{0.5\textwidth}
  \vspace{8pt}
  {\itshape \textbf{Fach:} .NET-Programmierung mit C\# 2026} \\
  {\itshape \textbf{Thema:} WeatherApp} \\
  {\itshape \textbf{Name:} Fadhil Fasalu Rahiman}
\end{minipage}
\begin{minipage}[t]{0.5\textwidth}
  \flushright
  \vspace{8pt}
  \itshape \textbf{Studiengang:} Computer Science \\
  \itshape \textbf{Datum:} \today
\end{minipage}

\section{Benutzerhandbuch}

\subsection{Beschreibung aller Funktionen}
WeatherApp ist eine Desktopanwendung zur schnellen Anzeige von
Wetterinformationen für frei suchbare Städte. Die Anwendung bietet
folgende Hauptfunktionen:

\begin{itemize}
  \item \textbf{Städtsuche:} Ermöglicht die Suche nach beliebigen
    Städten und lädt anschließend die zugehörigen Wetterdaten.
  \item \textbf{Aktuelle Wetteranzeige:} Zeigt die aktuelle
    Temperatur, den Wetterzustand sowie die Tageshöchst- und
    Tagestiefsttemperatur der ausgewählten Stadt an.
  \item \textbf{12-Stunden-Vorhersage:} Zeigt die Wetterentwicklung
    und Temperaturen für die kommenden zwölf Stunden ab der aktuellen
    Stunde.
  \item \textbf{7-Tage-Vorhersage:} Gibt einen Überblick über die
    Wetterbedingungen sowie die Höchst- und Tiefsttemperaturen der
    nächsten sieben Tage.
  \item \textbf{Favoriten:} Häufig verwendete Städte können
    gespeichert und anschließend schnell über die Favoritenliste
    aufgerufen werden.
  \item \textbf{Wetterempfehlungen:} Gibt abhängig von den aktuellen
    Wetterbedingungen praktische Hinweise, beispielsweise zu Regen,
    Kälte, Hitze oder starkem Wind.
  \item \textbf{Einstellungen:} Ermöglicht die Auswahl der
    Temperatureinheit (Celsius/Fahrenheit) und das Festlegen einer
    Standardstadt.
\end{itemize}

\subsection{Anleitung für unbedarfte Benutzer}
\begin{enumerate}
  \item \textbf{Start und Stadtsuche}
    \begin{itemize}
      \item Anwendung starten: Beim Start wird automatisch
        die gespeicherte Standardstadt geladen.
      \item Seitenleiste öffnen: Über das Menüsymbol kann
        die Seitenleiste geöffnet werden.
      \item Stadt suchen: Den Namen der gewünschten Stadt in
        das Suchfeld eingeben und auf "`Search"' klicken.
      \item Stadt auswählen: Eine Stadt aus der Ergebnisliste
        auswählen.
      \item Wetterdaten laden: Nach der Auswahl werden die
        Wetterdaten der Stadt im Hauptfenster angezeigt.
    \end{itemize}
  \item \textbf{Wetteranzeige}
    \begin{itemize}
      \item Oben werden der Stadtname, die aktuelle Temperatur und der
        aktuelle Wetterzustand angezeigt.
      \item Darunter befinden sich die Tageshöchst- und
        Tagestiefsttemperatur.
      \item Eine Wetterempfehlung gibt praktische Hinweise für den
        Alltag, beispielsweise zur Mitnahme eines Regenschirms oder
        zum Tragen warmer Kleidung.
      \item Die 12-Stunden-Vorhersage zeigt die Wetterentwicklung der
        kommenden Stunden ab der aktuellen Stunde.
      \item Die 7-Tage-Vorhersage gibt einen Überblick über die
        kommenden Tage mit Temperaturwerten und Wettersymbolen.
      \item Zusätzliche Detailkarten informieren über
        Windgeschwindigkeit und Niederschlag.
    \end{itemize}
  \item \textbf{Favoriten und Einstellungen}
    \begin{itemize}
      \item Über den ungefüllten Stern kann die aktuell ausgewählte
        Stadt zu den Favoriten hinzugefügt werden.
      \item Ein gefüllter Stern zeigt an, dass die Stadt bereits als
        Favorit gespeichert ist. Durch erneutes Anklicken wird sie aus
        den Favoriten entfernt.
      \item Gespeicherte Favoriten können über die Seitenleiste
        ausgewählt oder entfernt werden.
      \item In den Einstellungen kann zwischen Celsius und Fahrenheit
        als Temperatureinheit gewechselt werden.
      \item Die aktuell ausgewählte Stadt kann als Standardstadt
        gespeichert werden.
    \end{itemize}
  \item \textbf{Fehlerbehandlung}
    \begin{itemize}
      \item Bei API- oder Netzwerkfehlern wird eine verständliche
        Statusmeldung in der Anwendung angezeigt.
      \item Bei kritischen Fehlern während des Starts, beispielsweise
        bei Problemen mit der Datenbankinitialisierung oder
        -migration, erscheint ein entsprechendes Dialogfenster.
    \end{itemize}
\end{enumerate}

\section{Beschreibung der Anwendungsidee}
WeatherApp ist eine WPF-Desktopanwendung, die eine schnelle und
übersichtliche Darstellung von Wetterinformationen für frei wählbare
Städte ermöglicht. Neben aktuellen Wetterdaten und Vorhersagen bietet
die Anwendung praktische Wetterempfehlungen und ermöglicht es, häufig
verwendete Städte als Favoriten zu speichern. Ziel ist es, wichtige
Informationen für die tägliche Planung verständlich und ohne unnötige
Komplexität bereitzustellen.

\subsection{Zielgruppe}
Die Anwendung richtet sich an Personen, die regelmäßig die
Wetterbedingungen für ihre täglichen Aktivitäten, den Arbeitsweg oder
die Reiseplanung überprüfen möchten. Dazu gehören insbesondere
\textbf{Studierende, Pendler und Reisende}, die schnell und unkompliziert
Informationen über das Wetter an verschiedenen Orten benötigen.

\subsection{Most Valuable Features (MVF)}
\begin{itemize}
  \item \textbf{Aktuelles Wetter und Vorhersagen:} Zeigt die aktuellen
    Wetterbedingungen sowie eine 12-Stunden- und 7-Tage-Vorhersage für
    die ausgewählte Stadt an.
  \item \textbf{Städtesuche und Favoriten:} Ermöglicht die Suche nach
    Städten und das Speichern häufig verwendeter Orte für einen
    schnellen Zugriff.
  \item \textbf{Wetterempfehlungen:} Gibt auf Grundlage der aktuellen
    Wetterbedingungen praktische Hinweise, beispielsweise zur Mitnahme
    eines Regenschirms oder zum Tragen geeigneter Kleidung.
  \item \textbf{Temperatureinstellungen und Standardstadt:} Ermöglicht
    die Auswahl zwischen Celsius und Fahrenheit sowie das Festlegen
    einer bevorzugten Standardstadt.
\end{itemize}

\newpage
\section{Code-Statistik (Lines of Code)}
Die Ermittlung des Quellcode-Umfangs erfolgte strikt nach Vorgabe
über das Tool \texttt{cloc} für das Projekt
\texttt{WeatherApp} unter Ausschluss von Kompilierordnern (\texttt{bin},
\texttt{obj}), externen Bibliotheken und Test-Dateien.

Der Richtwert von mindestens 1500 Lines of Code (C\# + XAML) wird mit
insgesamt \textbf{2.064 Zeilen Code und Markup} deutlich erfüllt.

\vspace{0.8em}
\noindent
\renewcommand{\arraystretch}{1.2}
\begin{tabular}{|p{4cm}|r|r|r|r|}
  \hline
  \textbf{Sprache} & \textbf{Dateien} & \textbf{Leerzeilen} &
  \textbf{Kommentarzeilen} & \textbf{Codezeilen (LOC)} \\
  \hline
  XAML & 9 & 95 & 17 & 735 \\
  \hline
  C\# & 34 & 271 & 142 & 1.329 \\
  \hline
  \textbf{Gesamt (SUM)} & \textbf{43} & \textbf{366} & \textbf{159} &
  \textbf{2.064} \\
  \hline
\end{tabular}

\section{Verwendeter Fremdcode und externe Tools}
Das Projekt wurde als Windows Presentation Foundation (WPF)
Desktop-Anwendung in C\# unter .NET 10.0 entwickelt. Es nutzt die
folgenden externen NuGet-Pakete und Architektur-Tools:

\vspace{1em}
\noindent
\renewcommand{\arraystretch}{1.3}
\begin{tabular}{|p{4.5cm}|p{11.5cm}|}
  \hline
  \textbf{Paket / Tool} & \textbf{Zweck} \\
  \hline
  CommunityToolkit.\newline Mvvm & Unterstützt die Umsetzung des MVVM-Patterns,
  unter anderem durch ObservableObject, ObservableProperty und
  RelayCommand. \\
  \hline
  Entity Framework Core & ORM für den Zugriff auf die lokale
  SQLite-Datenbank und die Verwaltung von Datenbankmigrationen. \\
  \hline
  Microsoft.Data.Sqlite /\newline SQLitePCLRaw & Ermöglicht die Anbindung und
  Verwendung der SQLite-Datenbank. \\
  \hline
  Microsoft.Extensions.\newline DependencyInjection & Wird für Dependency
  Injection verwendet und ermöglicht die Verwaltung und Bereitstellung
  der Services und ViewModels. \\
  \hline
  Open-Meteo Forecast API & Liefert die Wetterdaten für die aktuelle
  Wetterlage sowie die Stunden- und Tagesvorhersagen. \\
  \hline
  Open-Meteo Geocoding API & Ermöglicht die Suche nach Städten und die
  Auflösung der entsprechenden Koordinaten. \\
  \hline
  xUnit & Framework für die Erstellung und Ausführung der
  automatisierten Unit-Tests. \\
  \hline
  Visual Studio / NuGet & Visual Studio dient als Entwicklungsumgebung;
  NuGet wird zur Verwaltung der verwendeten Pakete und Abhängigkeiten
  eingesetzt. \\
  \hline
  Git & Versionsverwaltung zur Nachverfolgung von Änderungen am
  Projekt. \\
  \hline
\end{tabular}

\newpage
\section{Einsatz von KI-Tools}
Im Rahmen dieser Studienarbeit wurden die KI-Assistenten \textbf{Claude}
und \textbf{Codex (ChatGPT)} unterstützend eingesetzt. Die KI wurde
dabei \textbf{nicht zur vollständigen Generierung der Anwendung}
verwendet, sondern gezielt zur Unterstützung bei einzelnen Teilaufgaben
und zur Verbesserung des Entwicklungsprozesses.

Die KI-Unterstützung umfasste insbesondere folgende Bereiche:

\begin{itemize}
  \item \textbf{Projektstruktur:} Unterstützung bei der Planung und
    Strukturierung des Projekts.
  \item \textbf{Debugging und Code Review:} Analyse von Fehlern sowie
    Überprüfung und Verbesserungsvorschläge für bestehenden Code.
  \item \textbf{XML-Dokumentation:} Unterstützung beim Erstellen und
    Formulieren von
    XML-Dokumentationskommentaren.
  \item \textbf{XAML-Styling und UI-Design:} Unterstützung bei der
    Gestaltung und Verbesserung der Benutzeroberfläche.
  \item \textbf{Unit-Tests:} Unterstützung bei der Erstellung und
    Erweiterung von Unit-Tests sowie bei der Identifikation geeigneter
    Testfälle.
  \item \textbf{Erklärungen und Refactoring:} Unterstützung beim
    Verständnis von Konzepten und bei der Entwicklung von Ideen zur
    Verbesserung und Umstrukturierung des Codes.
  \item \textbf{Dokumentation:} Unterstützung bei der Strukturierung,
    Formulierung und sprachlichen Überarbeitung der Projektdokumentation.
\end{itemize}

Sämtliche Vorschläge wurden kritisch geprüft, nachvollzogen und vor der
Integration an die bestehende Softwarearchitektur angepasst. Die
Implementierung und Überprüfung der Lösungen erfolgte eigenständig.
KI-Tools dienten dabei ausschließlich als unterstützende Werkzeuge im
Entwicklungsprozess.

\end{document}
